\section{One multiplicative prime--sign field}
\label{sec:tpc43-prime-sign}

Let \((\varepsilon_p)_{p\in\cP}\) be independent signs, uniformly
distributed on \(\{-1,1\}\), for every prime occurring in the finite
packet.  Extend the prime signs to the completely multiplicative phase
\(\varepsilon(n)=\prod_p\varepsilon_p^{v_p(n)}\), and define the
squarefree--supported random multiplicative function
\begin{equation}
 F_\varepsilon(n)
 :=\mu^2(n)\prod_{p\mid n}\varepsilon_p.
 \label{eq:tpc43-rademacher-multiplicative}
\end{equation}
Thus \(F_\varepsilon\) is multiplicative on coprime inputs and has the
same zero set and absolute values as \(\mu\).  It is not called
completely multiplicative because its value is zero on nonsquarefree
integers.  At the distinguished prime point
\begin{equation}
 \varepsilon_p=-1\quad(p\in\cP),
 \label{eq:tpc43-all-minus-point}
\end{equation}
one has
\begin{equation}
 F_\varepsilon(n)=\mu(n)
 \label{eq:tpc43-all-minus-is-mobius}
\end{equation}
for every relevant integer.

The deformation is applied consistently to the two M\"obius occurrences
in each physical row.  Define
\begin{equation}
 [\mathscr T_{k,\varepsilon}^L]_{\gamma,j}
 :=\sum_\alpha
 a^L_{\alpha,\gamma,j,k}
 F_\varepsilon(d_\alpha)
 F_\varepsilon(n_{\alpha,j,k}),
 \label{eq:tpc43-random-left-column}
\end{equation}
and define \(\mathscr T_{k,\varepsilon}^R\) symmetrically.  In
particular,
\begin{equation}
 \mathscr T_{k,-\one}^{L/R}=\mathscr T_k^{L/R}.
 \label{eq:tpc43-physical-corner}
\end{equation}
This is not an independent row randomization: the same prime sign appears
at every source and target integer divisible by that prime.

For a squarefree positive integer \(s\), write
\begin{equation}
 \chi_s(\varepsilon):=\prod_{p\mid s}\varepsilon_p,
 \qquad \chi_1=1.
 \label{eq:tpc43-Walsh-character}
\end{equation}
These are the Walsh characters of the finite prime cube and satisfy
\begin{equation}
 \E_\varepsilon
 \chi_s(\varepsilon)\chi_t(\varepsilon)
 =\one_{s=t}.
 \label{eq:tpc43-Walsh-orthogonality}
\end{equation}

If both \(d\) and \(n\) are squarefree, put
\begin{equation}
 g=(d,n),
 \qquad
 s(d,n):=\frac{dn}{g^2}.
 \label{eq:tpc43-squarefree-kernel}
\end{equation}
Then \(s(d,n)\) is squarefree and
\begin{equation}
 F_\varepsilon(d)F_\varepsilon(n)
 =\chi_{s(d,n)}(\varepsilon).
 \label{eq:tpc43-fused-Walsh-character}
\end{equation}
If \(n\) is not squarefree, both sides of every physical atom in which
it occurs are interpreted as zero.  Therefore no conditioning on
squarefreeness is hidden in the model.

\begin{remark}[Relation to external row phases]
\label{rem:tpc43-not-row-randomization}
Independent signs attached directly to row labels would diagonalize any
finite row family and would be a generic phase--frame statement of the
kind already delimited in TPC--40.  In
\eqref{eq:tpc43-random-left-column}, the signs live at primes, respect
multiplication, occur simultaneously at the source and target, and have
the actual M\"obius function as the point
\eqref{eq:tpc43-all-minus-point}.  The product--fiber geometry needed to
control their collisions is the arithmetic content of the next section.
\end{remark}
